\section{Proof Sketches}

We provide proof sketches emphasizing where each conclusion originates.
No additional assumptions beyond OC/ETS structure are used.

\subsection{Sketch for Theorem~\ref{thm:saturation}}

Because $\Omega$ is defined by threshold inequalities $\ThetaSys(\omega)\le 0$,
points near $\partial\Omega$ are \emph{constraint-tight}: small moves in ambient
space typically violate at least one component of $\ThetaSys$.
If $\mathcal{I}$ ``tries'' to enact a boundary-crossing change, it must either
(i) step outside admissibility (violating $\ThetaSys$) or (ii) refuse to act on
those states (domain restriction). This yields $\Omega^\ast$ as the set of
states where constraints are tight enough that a nontrivial boundary-crossing
map cannot be defined while preserving $\ThetaSys$.

\subsection{Sketch for Theorem~\ref{thm:no_free}}

The boundary is where admissibility is fragile: the feasible directions in state
space collapse as constraints tighten. A $\Theta$-preserving boundary-crossing
operator cannot move ``freely''; it must route along constraint-compatible
directions.

Two structural failure modes are unavoidable:
\begin{enumerate}
\item \emph{Continuity loss.} Constraint routing introduces fragmentation:
previously connected regions of $\Omega$ become disconnected under the mapping
(or trajectories concentrate into narrower subsets), reducing $k$ for some
states.
\item \emph{Phase transition.} If the operator enforces constraint satisfaction
by cutting off cycles/flows required for closure (a core OC notion), then the
system crosses a phase boundary into $\INERTIA$, $\COLLAPSE$, or $\STOP$.
\end{enumerate}

Thus, if an operator both crosses $\partial\Omega$ and preserves $\ThetaSys$,
it cannot also guarantee $\Delta k\ge 0$ and phase preservation everywhere on a
nontrivial domain.

\subsection{Sketch for Proposition~\ref{prop:trivial}}

If an operator never confronts $\partial\Omega$, it is confined to interior
movement within an already admissible region. In OC terms, it cannot realize
structural changes whose feasibility is precisely blocked by $\ThetaSys$ at the
boundary: it only re-parameterizes within existing degrees of freedom. Hence it
may be safe (non-decreasing $k$, phase-preserving) but is correspondingly
limited in reach.

\subsection{Sketch for Theorem~\ref{thm:stop}}

$\STOP$ is terminal in OC/ETS terms: it corresponds to a death condition such as
loss of admissible cycles, collapse of closure, or emptiness of feasible
continuation under thresholds. Any subsequent $\Theta$-preserving intervention
would need an admissible path back to $\OK$, but terminality rules out such a
path by definition of the phase partition and death threshold logic.
